\usepackage[T1]{fontenc}
\usepackage[a4paper,margin=26mm]{geometry}
\usepackage{amsmath,amssymb,amsthm,mathtools}
\usepackage{booktabs,longtable,tabularx,array,enumitem}
\usepackage{microtype,xurl}
\usepackage{tikz}
\usetikzlibrary{arrows.meta,positioning}
\usepackage[hidelinks]{hyperref}
\usepackage{listings}
\lstset{basicstyle=\ttfamily\small,columns=fullflexible,breaklines=true,
  keepspaces=true,showstringspaces=false,frame=single,framerule=.2pt}
\setlist{itemsep=3pt,topsep=5pt}
\newtheorem{theorem}{Theorem}[section]
\newtheorem{proposition}[theorem]{Proposition}
\newtheorem{lemma}[theorem]{Lemma}
\newtheorem{remark}[theorem]{Remark}
\newtheorem{exercise}[theorem]{Exercise}
\newcommand{\R}{\mathbb R}
\newcommand{\tr}{\operatorname{tr}}
\newcommand{\Id}{\operatorname{Id}}
\newcommand{\curl}{\operatorname{curl}}
\newcommand{\F}{\mathcal F}
\newcommand{\energy}[1]{\left\|#1\right\|_a}
\DeclareRobustCommand{\code}[1]{\texttt{\detokenize{#1}}}
\pdfstringdefDisableCommands{\def\code#1{\detokenize{#1}}}
% Snapshot of manuscript labels; regenerate with refresh_sources.py.
\expandafter\def\csname manuscript:eq:torsion-state\endcsname{1}
\expandafter\def\csname manuscript:eq:J\endcsname{2}
\expandafter\def\csname manuscript:eq:F\endcsname{3}
\expandafter\def\csname manuscript:thm:conditional\endcsname{2.1}
\expandafter\def\csname manuscript:eq:area\endcsname{4}
\expandafter\def\csname manuscript:eq:area-gradient\endcsname{5}
\expandafter\def\csname manuscript:eq:area-hessian\endcsname{6}
\expandafter\def\csname manuscript:eq:material\endcsname{7}
\expandafter\def\csname manuscript:eq:HJ\endcsname{8}
\expandafter\def\csname manuscript:eq:HF-general\endcsname{9}
\expandafter\def\csname manuscript:eq:critical-gradient\endcsname{10}
\expandafter\def\csname manuscript:eq:HF-critical\endcsname{11}
\expandafter\def\csname manuscript:eq:additive\endcsname{12}
\expandafter\def\csname manuscript:sec:fan-explicit\endcsname{3.2}
\expandafter\def\csname manuscript:eq:fan-hat-gradient\endcsname{14}
\expandafter\def\csname manuscript:eq:dihedral-covariance\endcsname{15}
\expandafter\def\csname manuscript:eq:XYZQ\endcsname{16}
\expandafter\def\csname manuscript:eq:sector-identities\endcsname{17}
\expandafter\def\csname manuscript:prop:XYZ-anisotropy\endcsname{3.1}
\expandafter\def\csname manuscript:eq:sector-anisotropy\endcsname{19}
\expandafter\def\csname manuscript:eq:XYZ-delta\endcsname{20}
\expandafter\def\csname manuscript:eq:delta-elementary-bound\endcsname{21}
\expandafter\def\csname manuscript:eq:mixed-cell-scaling\endcsname{24}
\expandafter\def\csname manuscript:eq:mixed-cell-Hadamard\endcsname{25}
\expandafter\def\csname manuscript:eq:XYZ-mixed-cell\endcsname{26}
\expandafter\def\csname manuscript:ex:XYZ-triangle\endcsname{3.2}
\expandafter\def\csname manuscript:eq:XYZ-triangle\endcsname{28}
\expandafter\def\csname manuscript:lem:local-blocks\endcsname{3.3}
\expandafter\def\csname manuscript:eq:local-three-blocks\endcsname{30}
\expandafter\def\csname manuscript:eq:area-symbol-explicit\endcsname{31}
\expandafter\def\csname manuscript:eq:local-symbol-expanded\endcsname{32}
\expandafter\def\csname manuscript:eq:local-symbol-entries\endcsname{33}
\expandafter\def\csname manuscript:eq:explicit-cancellation\endcsname{34}
\expandafter\def\csname manuscript:eq:Hrad\endcsname{36}
\expandafter\def\csname manuscript:prop:block-circ\endcsname{4.1}
\expandafter\def\csname manuscript:eq:block-circ\endcsname{37}
\expandafter\def\csname manuscript:eq:symbol\endcsname{38}
\expandafter\def\csname manuscript:eq:symbol-entries\endcsname{39}
\expandafter\def\csname manuscript:eq:mode-eigs\endcsname{40}
\expandafter\def\csname manuscript:prop:modes\endcsname{4.2}
\expandafter\def\csname manuscript:eq:G-three\endcsname{42}
\expandafter\def\csname manuscript:prop:explicit-symbol\endcsname{4.3}
\expandafter\def\csname manuscript:eq:alpha-explicit\endcsname{44}
\expandafter\def\csname manuscript:eq:beta-explicit\endcsname{45}
\expandafter\def\csname manuscript:eq:gamma-explicit\endcsname{46}
\expandafter\def\csname manuscript:eq:trace-no-delta\endcsname{47}
\expandafter\def\csname manuscript:eq:mixed-cell-penalties\endcsname{49}
\expandafter\def\csname manuscript:eq:mode-zero-gram\endcsname{50}
\expandafter\def\csname manuscript:eq:translation-symbol-identities\endcsname{51}
\expandafter\def\csname manuscript:eq:conjugate-symbols\endcsname{52}
\expandafter\def\csname manuscript:sec:validated\endcsname{5}
\expandafter\def\csname manuscript:sec:certificate-target\endcsname{5.1}
\expandafter\def\csname manuscript:prop:real-fourier-pairings\endcsname{5.1}
\expandafter\def\csname manuscript:eq:real-pairing-symbol\endcsname{53}
\expandafter\def\csname manuscript:sec:residual-principle\endcsname{5.2}
\expandafter\def\csname manuscript:eq:equilibrated-majorant\endcsname{54}
\expandafter\def\csname manuscript:eq:curl-state-flux\endcsname{55}
\expandafter\def\csname manuscript:sec:hessian-residual\endcsname{5.3}
\expandafter\def\csname manuscript:eq:pullback-map\endcsname{56}
\expandafter\def\csname manuscript:eq:pulled-forms\endcsname{57}
\expandafter\def\csname manuscript:eq:pullback-coefficients\endcsname{58}
\expandafter\def\csname manuscript:eq:d-pullback\endcsname{59}
\expandafter\def\csname manuscript:eq:discrete-differentiated-systems\endcsname{60}
\expandafter\def\csname manuscript:eq:Cqr\endcsname{61}
\expandafter\def\csname manuscript:eq:second-lifting-functional\endcsname{62}
\expandafter\def\csname manuscript:eq:second-lifting\endcsname{63}
\expandafter\def\csname manuscript:prop:second-residual\endcsname{5.2}
\expandafter\def\csname manuscript:eq:second-residual-bound\endcsname{64}
\expandafter\def\csname manuscript:eq:energy-error-identity\endcsname{65}
\expandafter\def\csname manuscript:eq:second-residual-identity\endcsname{66}
\expandafter\def\csname manuscript:sec:algebraic-errors\endcsname{5.4}
\expandafter\def\csname manuscript:eq:algebraic-residual\endcsname{67}
\expandafter\def\csname manuscript:eq:four-algebraic-residuals\endcsname{68}
\expandafter\def\csname manuscript:eq:algebraic-transfer\endcsname{69}
\expandafter\def\csname manuscript:prop:corrected-centers\endcsname{5.4}
\expandafter\def\csname manuscript:eq:center-identity\endcsname{70}
\expandafter\def\csname manuscript:eq:J-center-identity\endcsname{71}
\expandafter\def\csname manuscript:eq:center-algebraic-bound\endcsname{72}
\expandafter\def\csname manuscript:sec:entry-enclosure\endcsname{5.5}
\expandafter\def\csname manuscript:eq:combined-deltas\endcsname{73}
\expandafter\def\csname manuscript:eq:combined-BJ\endcsname{74}
\expandafter\def\csname manuscript:eq:corrected-F-center\endcsname{75}
\expandafter\def\csname manuscript:eq:complete-F-radius\endcsname{76}
\expandafter\def\csname manuscript:sec:certificate-implementation\endcsname{5.6}
\expandafter\def\csname manuscript:sec:rate\endcsname{6}
\expandafter\def\csname manuscript:eq:w-rates\endcsname{77}
\expandafter\def\csname manuscript:eq:U-rate\endcsname{78}
\expandafter\def\csname manuscript:prop:h2\endcsname{6.1}
\expandafter\def\csname manuscript:conj:h2\endcsname{6.2}
\expandafter\def\csname manuscript:eq:h2-conj\endcsname{80}
\expandafter\def\csname manuscript:tab:rates\endcsname{1}
\expandafter\def\csname manuscript:eq:second-bound-rate-data\endcsname{82}
\expandafter\def\csname manuscript:sec:pentagon\endcsname{7}
\expandafter\def\csname manuscript:eq:n5-eigs\endcsname{84}
\expandafter\def\csname manuscript:tab:pentagon-entry-radii\endcsname{2}
\expandafter\def\csname manuscript:eq:largest-entry-certificate\endcsname{85}
\expandafter\def\csname manuscript:eq:validated-count\endcsname{86}
\expandafter\def\csname manuscript:thm:pentagon\endcsname{7.1}
\expandafter\def\csname manuscript:sec:n3-n10\endcsname{8}
\expandafter\def\csname manuscript:tab:n3-n10-certificates\endcsname{3}
\expandafter\def\csname manuscript:tab:n3-n10-margins\endcsname{4}
\expandafter\def\csname manuscript:tab:m64-certificates\endcsname{5}
\expandafter\def\csname manuscript:tab:m128-certificates\endcsname{6}
\expandafter\def\csname manuscript:thm:n3-n10\endcsname{8.1}
\expandafter\def\csname manuscript:sec:reproducibility\endcsname{9}
\expandafter\def\csname manuscript:app:certificate-details\endcsname{A}
\expandafter\def\csname manuscript:app:pullback\endcsname{A.1}
\expandafter\def\csname manuscript:eq:Aq-pullback\endcsname{87}
\expandafter\def\csname manuscript:eq:Aqr-pullback\endcsname{88}
\expandafter\def\csname manuscript:app:assembly\endcsname{A.2}
\expandafter\def\csname manuscript:eq:alpha-lower\endcsname{89}
\expandafter\def\csname manuscript:eq:square-lambda-bound\endcsname{90}
\expandafter\def\csname manuscript:eq:analytic-local-matrices\endcsname{91}
\expandafter\def\csname manuscript:eq:analytic-local-stiffness\endcsname{92}
\expandafter\def\csname manuscript:eq:analytic-differentiated-matrices\endcsname{93}
\expandafter\def\csname manuscript:app:fluxes\endcsname{A.3}
\expandafter\def\csname manuscript:eq:second-source-divergence\endcsname{94}
\expandafter\def\csname manuscript:eq:flux-majorant\endcsname{96}
\expandafter\def\csname manuscript:eq:G-material\endcsname{97}
\expandafter\def\csname manuscript:eq:curl-material-flux\endcsname{99}
\expandafter\def\csname manuscript:eq:material-majorant\endcsname{100}
\expandafter\def\csname manuscript:app:direct-bounds\endcsname{A.4}
\expandafter\def\csname manuscript:eq:gram-error\endcsname{102}
\expandafter\def\csname manuscript:eq:wquad-error\endcsname{103}
\expandafter\def\csname manuscript:eq:arb-majorants\endcsname{104}

\newcommand{\meq}[1]{\csname manuscript:#1\endcsname}
\setlength{\emergencystretch}{2em}
\hypersetup{pdfauthor={Course notes based on the torsion manuscript},
 pdfsubject={Shape Hessians, finite element error bounds, and validated computing}}
